\section{Problem statement and requirements}
\label{sec:problem}

We consider a production line supervised by an information platform that
provides data acquisition from plant nodes and controlled writes to plant
nodes, in the sense of the enterprise--control integration models of
ISA-95 \cite{isa95_2025} and batch control models of ISA-88
\cite{isa88_2010}; the line--product--recipe--run tuples used throughout
this paper correspond to ISA-88 control-recipe and batch identities. LLM
agent teams run inside the same platform. Each team
is led by a coordinator agent that decomposes a stated production goal into
worker tasks; agents perceive the plant through queries against acquired
telemetry and act through platform operations. We do not consider
real-time control loops: the platform occupies a supervisory position in
which decisions occur on second-to-minute cadences, and equipment-level
protection remains with the plant's own basic process control and safety
instrumented systems \cite{iec61511}. The platform is not a safety
barrier and makes no claim inside a safety-instrumented function.

\subsection{Failure model}
\label{sec:problem-fm}

Four failure classes motivate the design. They are stated from the plant
operator's perspective: each is an event that can occur, that operators must
afterwards reconstruct, and that current agent frameworks make difficult to
reconstruct.

\textbf{M1: Unconstrained machine-originated write.} An agent issues a
setpoint write that violates plant engineering context: outside the
node's safe range, outside the recipe window of the product currently in
production, or in conflict with another proposal in flight. Session-level
tool permissions cannot see this context: the write tool is either allowed
(and the violation passes) or denied (and no legitimate write is possible).

\textbf{M2: Unattributable decision.} A write is blocked (or executes)
and the operator cannot determine which control decided. If scope checks,
engineering interlocks, and approval requirements are evaluated in one
amalgamated policy, the audit log records an outcome but not its cause.
Attribution matters operationally: a block caused by an expired approval
window is fixed differently from one caused by a recipe window.

\textbf{M3: Stale or replayed approval.} A human approves a proposed write;
plant state changes before the write executes (a different recipe starts, a
batch closes, limits are revised); the write executes under the old
approval. The approval was real, but it no longer describes the operation
performed.

\textbf{M4: Unverified receipt.} The platform reports success because a
protocol stack acknowledged a transmission. The device may not have applied
the value. A readback would have shown so. Operators who rely on the
transport acknowledgment form beliefs about plant state that the device does
not support.

Prompt injection sharpens all four classes: an agent that reads plant-facing
content can be steered to request writes it was never asked for
\cite{greshake2023injection,liu2024pinjection}. The path must therefore
treat the \emph{requester} as untrusted even when the request is
well-formed.

\subsection{Requirements}
\label{sec:problem-req}

From the failure model we derive six requirements for a governed write path.

\textbf{R1 (Context-encoded scope).} Agent authority over writes must be a
configuration fact (bound to specific nodes, expressible per point, per
recipe, or per line) and must derive from plant structure (lines,
products, recipes, runs), not from prompt conventions.

\textbf{R2 (Layered admission).} Every write must cross an admission
interval formed by independent layers: a structural hard range at the node,
then parameter, product, and recipe limits. Layers must be evaluated in a
fixed order and the first excluding layer must be identifiable.

\textbf{R3 (Defense in depth).} No single layer may be load-bearing for
safety. If one layer is disabled or misconfigured, remaining layers must
still bound writes, and the hard node range must be bypassable by no
configuration flag.

\textbf{R4 (Bound, rechecked approval).} When a binding requires human
approval, the approval must record the proposed values, the applicable
limits, an expiry, and the active line--product--recipe--run tuple; the
dispatching stage must recheck the record against current state before
executing, so that M3 fails closed.

\textbf{R5 (Honest receipts).} Write outcomes must be graded into at least
three classes (readback-verified, transport-acknowledged, and unverified),
and the platform must not represent a transport acknowledgment as
device-side confirmation.

\textbf{R6 (Total audit).} Every admission decision, including every
rejection with its deciding layer, must be recorded in an append-only
ledger, so that M2 fails closed: an absent record is itself evidence.

R1--R2 address M1; R2's layer ordering plus R6 address M2; R4 addresses M3;
R5 addresses M4; R3 constrains the whole chain. Section~\ref{sec:arch}
presents the architecture that realizes these requirements, and
Section~\ref{sec:eval} evaluates them, including a four-arm ablation
designed to attribute interception to specific layers (R2--R3) and to
measure false blocks, the cost side of governance.
